\chapter{Conclusion and Future Work}\label{ch:conclusion}

\section{Summary of Findings}\label{sec:summary}

We set out to test whether a very small restoration network could split an
image into frequency bands and route each band, pixel by pixel, to
convolutional experts of different kernel sizes, and whether the learned
routing would be meaningful. We return to the four objectives of
Section~\ref{sec:objectives}.

\paragraph{O1, build and verify.} \fcsg{} has 97,088 parameters and costs
9.39 reported GFLOPs at $256\times256$, well inside the limits of five million
and ten. It passed every verification check before training: exact band
reconstruction, consistent frequency units across image sizes, gradient flow,
mixed-precision execution, single-image overfitting and exact checkpoint
restoration (Section~\ref{sec:verification}).

\paragraph{O2, trustworthy baselines.} Our FFDNet implementation reproduces
the published CBSD68 result to within 0.01\dB{} using the author's weights,
and our DnCNN training loop reaches 32.65\dB{} on DIV2K for Gaussian
denoising (Section~\ref{sec:baseline-repro}).

\paragraph{O3, matched comparison.} Under one composite degradation and one
training protocol, \fcsg{} reaches 27.24\dB{} PSNR, 0.753 SSIM and 0.396
LPIPS on 100 DIV2K validation images. This is 0.59\dB{} below DnCNN and
0.95\dB{} below blind FFDNet, and it loses to FFDNet on every one of the 100
images. It does so with 82.6\% fewer parameters than DnCNN and 88.6\% fewer
than FFDNet. Its low FLOP count, however, did not make it fast: on a Tesla T4
it is the slowest of the four models (Chapter~\ref{ch:results}).

\paragraph{O4, the routing.} The learned gate varies across the image,
selects the $3\times3$ expert almost everywhere in the high band, and shifts
its weights with the noise level, with within-image correlations up to
$r=+0.81$. Overriding its selections at inference costs 0.37\dB{} on average,
while flattening its weights within the selected pair costs only 0.04\dB{},
so the trained model depends on which experts are chosen more than on how
they are mixed (Chapter~\ref{ch:routing}).

In short, frequency-band routing produced a small model whose behaviour can
be inspected and related to the degradation, but within our budget it did
not produce a model that is better or faster than simple convolutional
baselines.

\section{Limitations}\label{sec:limitations}

Several limitations bound how far these results can be taken.

Each model was trained once with an unfixed seed, so we have no estimate of
run-to-run variance. The consistency of the gap across images makes it
unlikely to be an accident of one seed, but small differences, such as
between the two FFDNet variants, cannot be interpreted. All results are on
the DIV2K validation split, which was also used to monitor training; there is
no untouched test set.

The degradation is synthetic. Real archival photographs suffer from fading,
scratches, dust, colour casts and scanner artefacts that our pipeline does
not model, and we did not evaluate on real archival images or ask people to
judge the results. The good performance of any of these models on our data
says nothing yet about real archives.

The routing analysis uses ten source images. Its correlations are
exploratory, the four degradation parameters vary together, and there is no
correction for the number of comparisons. The inference interventions
measure the trained model's dependence on its gate, not the value of having a
gate during training.

The cost measurements cover one GPU, the Tesla T4, at batch size 1 in FP32.
We did not profile individual operators, so our explanation of why \fcsg{} is
slow is a reasoned account rather than a measurement. We did not measure
latency on CPUs or mobile hardware, nor peak training memory at batch 16.

Finally, \fcsg{} computes all nine expert passes regardless of the gate, so
its sparse routing does not save computation.

\section{Future Work}\label{sec:future}

The most important next step is a matched training ablation: train \fcsg{}
and an otherwise identical model with fixed uniform weights from the same
initialisation and seed, and compare them on the full validation set,
preferably with several seeds each. This is the only experiment that can
show whether learning a router helps, rather than whether a trained model
depends on its router. Related ablations would replace the three kernel sizes
with three identical $3\times3$ experts, replace the spatial gate with a
single per-image gate, replace SE fusion with plain concatenation, and
replace the soft masks with hard ones. Each of these isolates one design
decision of Chapter~\ref{ch:method}.

The second direction is to turn sparse weights into sparse computation. Since
the trained model mostly cares about which two experts are chosen, skipping
the third expert in tiles where it is never selected could reduce the
running cost by up to a third. This should be preceded by an operator-level
profile, because if latency is dominated by memory traffic and launch
overhead, removing arithmetic may not help much. Measuring latency on a CPU
and on a mobile accelerator would show whether the model's low FLOP count
is worth more on hardware that suits depthwise convolutions.

Third, the capacity gap could be addressed directly by widening the experts
within the parameter budget, which is about fifty times larger than what the
model uses. The compute budget is nearly exhausted, so this would have to go
together with cheaper expert evaluation, for example by sharing one
unshuffled feature map across bands. Letting the gate see a larger context than a $16\times16$
window may also help.

Finally, any claim about archival restoration needs real archival
photographs, a controlled degradation sweep to separate the effect of each
corruption on the routing, and a held-out test set that was never used
during development.
